\chapter{结果与讨论}

\section{候选事例与拟合结果}

本节用于汇总最终信号候选数、背景估计、拟合结果和显著性。
若后续分析给出的是观测结果，则应明确统计显著性和对应检验方法；
若给出的是上限结果，则应明确置信区间构造方式及其输入假设。

当前应保留的结果占位包括：

\begin{itemize}
  \item \textbf{[待填：最终候选事例数]}
  \item \textbf{[待填：拟合信号产额]}
  \item \textbf{[待填：统计显著性或上限结果]}
\end{itemize}

\section{修正后结果与物理解释}

如果后续输出包含截面、比值或微分分布，本节应说明其定义、误差组成和适用相空间。
如果当前阶段只完成候选观测或方法学验证，本节则应聚焦于结果含义、限制条件和下一步工作。

\section{与相关工作的比较}

这里可以结合 \texttt{struct-ref/} 中公开结果的组织方式，讨论本文结果与已有双 quarkonium、
三 quarkonium 或相关伴随产生研究之间的联系。
但在没有最终数值和可靠可比定义之前，不应写出定量优劣判断。
